\chapter{Algebraic sets}
\section{Basic definitions}
Some important observations:
\begin{itemize}
    \item If $\bbF$ is a finite field, the map which takes a polynomial in $\bbF[X_1,\dots,X_n]$ and 
    send to a function $\bbF^n\to \bbF$ is not an isomorphism. For example, if $\bbF$ 
    has $q$ elements, then $\prod_{q\in\bbF}(X-q)\in\bbF[X]\mapsto 0$. However, if $\bbF$ is infinity and 
    $a_nX^n + a_{n-1}X^{n-1}+\cdots + a_0 = 0$, then we can select distinc elements $\alpha_0,\dots,\alpha_{n}$ and 
    we conclude that 
    \begin{equation}
        \begin{pmatrix}a_0&\dots&a_n\end{pmatrix}
       \begin{pmatrix}
            1&1&\cdots&1\\\alpha_0&\alpha_1&\cdots&\alpha_{n}\\\vdots&\vdots&\cdots&\vdots\\
            \alpha_0^n&\alpha_1^n&\cdots&\alpha_{n}^{n}
        \end{pmatrix}=0.
    \end{equation}
    However $\det(\alpha_i^j)_{i,j=0}^n=\prod_{i<j}(\alpha_i-\alpha_j)\neq0$ and we conclude that $a_0,\dots,a_n=0$. 
    Now, if $f\in \bbF[X_1,\dots,X_n]$, then we can write (assumin $X_1$ appears in f):
    \begin{equation}
        f = g_d(X_2,\dots,X_n)X_1^d +\cdots + g_1(X_2,\dots,X_n)X_1 + g_0(X_2,\dots,X_n).
    \end{equation}
    So the equality $f=0$, the previous step for the case $\bbF[X]$ and an induction on the number of variables yields
    $g_d,\dots,g_0 =0$. Then we see that if $|\bbF|=\infty$, we can identify the ring $\bbF[X_1,\dots,X_n]$ with 
    polynomial functions $\bbF^n\to\bbF$.\par 

    As a consequence of this story, we see that if $|\bbF|=\infty$, then 
    \begin{equation}
        I(\bbF^n)=(0)\subset\bbF[X_1,\dots,X_n].
    \end{equation}

    \item(Hilbert basis theorem). The Hilbert basis theorem states that if $A$ is noetherian, then $A[X]$ is noetherian. 
    Since $\bbF[X]$ is a PID (and in consequence noetherian), we conclude by an induction on the number of variables that 
    $\bbF[X_1,\dots, X_n]$ is noetherian. 

    \begin{remark}
        A ring is said to be noetherian if it satisfy the ascending chain condition, or equivalently, 
         if all of its ideals are finitely generated.
    \end{remark}
    \item(Algebraic sets). Let $S \subset \bbF[X_1,\dots, X_n]$, then $V(S) = \{a\in\bbF^n: f(a)=0\}$. It's 
    immediate that $V(S) = V((S))$. The relation $V(\bigcup_{i\in I}S_i)=\bigcap_{i\in I}V(S_i)$ allow us 
    to define the Zariski topology on $\bbF^n$. 
    Moreover, for any subset $V\subset \bbF^n$, we let 
    $I(V)=\{f\in\bbF[X_1,\dots,X_n]: f(a)=0,\;\forall a\in V\}$. Then we have 
    \begin{equation}
       \left\{ \begin{aligned}
            &VI(W) = \overline{W}\\
            &IV(\fr{a}) = \sqrt{\fr{a}}
        \end{aligned}\right.
    \end{equation}
    We also have the important relation
    \begin{equation}
        V(\fr{a}\cap\fr{b}) = V(\fr{a}\fr{b}) = V(\fr{a})\cup V(\fr{b}).
    \end{equation}
    In fact 
    \begin{equation}
        V(\fr{a})\cup V(\fr{b})\underset{\fr{a}\cap\fr{b}\subset \fr{a},\fr{b}}{\subset}
        V(\fr{a}\cap\fr{b})\underset{\fr{a}\fr{b}\subset\fr{a}\cap\fr{b}}{\subset}V(\fr{a}\fr{b})
        =V(\{fg:f\in\fr{a},g\in\fr{b}\})\subset V(\fr{a})\cup V(\fr{b}).
    \end{equation}

    \item(Coordinate rings). Let $\fr{a} = I(V)$ where $V$ is an algebraic set. We define the coordinate ring 
    $A[V] = \bbF[X_1,\dots, X_n]/\fr{a}$.

    \item(LEX order). We define an ordering on $\bbZ_{\geq0}^n$ by $\alpha > \beta$ if 
    for the smallest $i$ such that $\alpha_i\neq \beta_i$ we have $\alpha_i > \beta_i$. So, with a 
    fixed LEX on $\bbF[X_1,\dots, X_n]$, we can definte the leading term of a polynomial: 
    $\text{LT}(f) = a_\alpha X^\alpha$, where $\alpha$ is the smallest tuple, relative to LEX such that 
    $a_\alpha\neq0$. In the case $\LT(f) = a_\alpha X^\alpha$, we also write $\LM(f) =X^\alpha$.

    \item(Groebner Basis). Let $\fr{a}$ be an ideal in $\bbF[X_1,\dots, X_n]$. A groebner basis of $\fr{a}$ 
    is a finite subset $G\subset\fr{a}$ such that $(\LT(g):g\in G) = \LT(\fr{a})$.\par 
    
    \begin{example}
        Let $R=\bbF[X_1,\dots, X_n, Y_1,\dots, Y_n]$ with LEX: $X_1>\cdots>X_n>Y_1>\cdots>Y_n$ and define an ideal 
    $\fr{a} = (X_1 - g_1,\dots, X_n - g_n)$ for $g_i\in\bbF[Y_1,\dots, Y_n]$. Since 
    $\LT(X_i - g_i)=X_i$, it's clear that $(X_1,\dots, X_n)\subset \LT(I)$. Now, if $f\in R$ and 
    $X_j\nmid \LT(f),\;\forall j$, then $f\in \bbF[Y_1,\dots, Y_n]$ because its smallest term has 
    power $\alpha \in \{0\}\times\bbZ^n$. So, we need to show that $I\cap \bbF[Y_1,\dots, Y_n] = (0)$. This can be seen 
    by proving that the kernel of $\varphi: R\to \bbF[Y]:X_i\mapsto g_i, Y_i\mapsto Y_i$ is $I$. The only nontrivial inclusion 
    is $\Ker(\varphi)\subset I$. To see this use a division algorithm to write 
    \begin{equation}
        f = h_1(X_1-g_1)+\cdots + h_n(X_n-g_n) + r
    \end{equation}
     where $r\neq0$ implies that no term of $r$ is divisible by $\LT(X_i - g_i) = X_i,\;\forall i$. By applying 
     $\varphi$ on both sides of this equation we have the desired inclusion. From this we see that 
     $G = (X_1-g_1,\dots, X_n-g_n)$ is a groebner basis for $I$.
    \end{example}
    \item(Hilbert Polynomial). 
    \begin{definition}
        Let $I\subset \bbF[X_1,\dots, X_n]$ be an ideal. We deffine the affine Hilbert function:
        \begin{equation}
            \upleft{a}\text{HF}_I:\bbZ_{>0}\to\bbZ_{>0}:s\mapsto \dim_{\bbF}(\frac{\bbF[X_1,\dots,X_n]_{\leq s}}{I_{\leq s}}).
        \end{equation}
    \end{definition}

    \begin{example}
        For example, the number $\upleft{a}\text{HF}_{(0)}(s)$ is the number of $n$-tuples of 
        nonnegative integers whose sum is $\leq s$. To compute this, write the generating function
        \begin{equation}
            \prod_{i=1}^n\frac{1}{1-X_i} = \sum_{\alpha\in\bbZ_{\geq0}^n}X^\alpha.
        \end{equation}
        Now put $X_i = t,\;\forall i$. We see that 
        \begin{equation}
            (1-t)^{-n} = \sum_{s\geq0}\binom{-n}{s}(-1)^st^s = \sum_{s\geq0}\binom{n+s}{s}t^s.
        \end{equation}
        So $\upleft{a}\text{HF}_{(0)}(s)=\binom{n+s}{s}=\frac{1}{n!}s^n + \cdots$.
    \end{example}
    \begin{example}
        Let $f\in\bbF[X_1,\dots, X_n]$ be a degree $d$ polynomial. Then 
        \begin{equation}
            \upleft{a}\text{HF}_{(f)}(s) = \binom{s+n}{s} - \binom{s-d+n}{s-d} = \frac{d}{(n-1)!}s^{n-1}+\cdots.
        \end{equation} 
    \end{example}
    An important property of the affine Hilbert function is the following:
\begin{theorem}
    Let $I\subset \bbF[X_1,\dots, X_n]$ be an ideal, then 
    \begin{equation}
        \upleft{a}\text{HF}_I(s)  = \upleft{a}\text{HF}_{\LT(I)}(s),\;s\geq0.
    \end{equation}
\end{theorem}
Another main result is the existence of the Hilbert polynomial.
\begin{theorem}
    There exist a polynomial $\upleft{a}\text{HP}_I(t)\in\bbQ[t]$ such that 
    \begin{equation}
        \upleft{a}\text{HP}_I(s) =   \upleft{a}\text{HF}_I(s),\;\forall s\geq0.
    \end{equation}
\end{theorem}
Proof (sketch):
\begin{itemize}
    \item Find a Groebner basis $G$ of $I$
    \item Let $M\subset \bbZ_{\geq0^n}$ be such that $\{X^\beta\}_{\beta\in M} = \{\LM(g):g\in G\}$
    \item Write 
    \begin{equation}
        C(I) = \bigcap_{\beta\in M}C(\beta),C(\beta)=\{\alpha |X^\beta\nmid X^\alpha\}
    \end{equation}
    as a union of $C(J,\tau)=\{\alpha|\alpha_j=\tau(j)\;\forall j\in J\}$. To do this use 
    \begin{equation}
        C(\beta) = \bigcup_{i=1}^n\bigcup_{t_i=0}^{\beta_i-1}C(\{i\},\;\tau:i\mapsto t_i).
    \end{equation}
    \item Use the formula $|C(J,\tau)_{\leq s}| = \binom{d+s-|\tau|}{d}, d = n-|J|, 
    |\tau| = \sum_{j\in J}\tau(j)$ and the inclusion-exclusion principle to compute 
    \begin{equation}
       \upleft{a}\text{HP}_I(s) = |C(I)_{\leq s}| = \biggl|\bigcup_{(J,\tau)}C(J,\tau)_{\leq s}\biggr|.
    \end{equation} 
\end{itemize} 
\begin{example}
    Let $I = (X_1-X_3^2,X_2-X_3^3)$ and fix the LEX order with $X_1>X_2>X_3$. Then we calculate
    \begin{itemize}
        \item $G=\{X_1-X_3^2,X_2-X_3^3\}$
        \item $M=\{(1,0,0),(0,1,0)\}$
        \item We have 
        \begin{equation}
            \begin{split}
                C(I) &= C(\{1\},\tau:1\mapsto0)\cap C(\{2\},\tau:2\mapsto0)\\
                &= C(\{1,2\}, \tau:1\mapsto 0, 2\mapsto0).
            \end{split}
        \end{equation}
    \end{itemize}
     So we conclude that $|C(I)|_{\leq s} = \binom{1+s-0}{1}=s+1$.
\end{example}
\item(Algebraic dimension). Here we define the dimension of an algebraic set.
\begin{definition}
    Let $V\subset \bbF^n$ be an algebraic set. We define its dimension as $\dim(V) = \deg \upleft{a}\text{HP}_{I(V)}(t)$.
    It has the property that it's the biggest integer $d\geq1$ such that there exist indexes 
    $1\leq i_1<\dots< i_d\leq n$ such that $I(V)\cap\bbF[X_{i_1},\dots, X_{i_d}]=(0)$. 
\end{definition}

\begin{example}
    Let $H_f = V(f)$, $f$ irreducible, then 
    \begin{equation}
        \begin{split}
            \dim(H_f) & = \deg\upleft{a}\text{HP}_{I(H_f)}(t)\\
            &=\deg \upleft{a}\text{HP}_{(f)}(t)\\
            &=\deg (\frac{\deg(f)}{(n-1)!}t^{n-1}+\cdots)\\
            &=n-1.
        \end{split}
    \end{equation}
\end{example}
The dimension has the expected property: If $W\subset V$ is closed, then $\dim(W)\leq\dim(V)$. In fact
\begin{equation}
    \begin{split}
        &I(V)\subset I(W)\Rightarrow\\
        &I(V)_{\leq s}\subset I(W)_{\leq s}\Rightarrow \\
        &\upleft{a}\text{HF}_{I(W)}(s) \leq \upleft{a}\text{HF}_{I(V)}(s),\;\forall s\geq0\Rightarrow\\
        &\deg \upleft{a}\text{HF}_{I(W)}(t)\leq \deg \upleft{a}\text{HF}_{I(V)}(t).
    \end{split}
\end{equation}
\item(Transcendence degree). Let $A$ be a field which is also a finitely generated $\bbF$-algebra. We define 
the transcendence degree of $A$ over $\bbF$ (and write $\partial_{\bbF}(A)$) to be largest size of an algebraicaly independent subset of $A$ over 
$\bbF$. 
\begin{example}
    We have $\partial_{\bbF}(\bbF[X_1,\dots, X_n])=n$. Indeed, let $\{f_1,\dots, f_r\}$ be polynomials and
    suppose that $r>n$. 
    Write 
    \begin{equation}
        \begin{split}
            &W_s = \{f^\beta:\beta\in\bbZ_{\geq0}^r,\;|\beta|\leq s\},\;d=\max\{\deg(f_i):i=1,\dots,r\}\\
            &V_s=\{X^\alpha:|\alpha|\leq s\}
        \end{split}
    \end{equation}
    Then $W_s\subset V_{ds}$ and 
    \begin{equation}
        |W_s| = \binom{r+s}{s} \overset{s\gg0}{>}\binom{n+sd}{sd} = \dim(V_{sd})\geq \dim(W_s).
    \end{equation}
    So, for $s\gg0$ there exist $\{a_\beta\}_\beta\subset\bbF$ such that 
    $\sum_{\beta}a_\beta f^\beta=0$ and $\{f_1,\dots, f_r\}$ is algebraicaly dependent.
    \end{example}
    \begin{theorem}
        Let $V$ be an algebraic set and write $I= I(V)$, then 
        \begin{equation}
            \dim(V) = \partial_{\bbF}(A[V])=\partial_{\bbF}(\frac{\bbF[X_1,\dots, X_n]}{I}).
        \end{equation}
    \end{theorem}
    Proof: Let $d=\dim(V)$, so there exist indexes $1\leq i_1<\cdots<i_d\leq n$ such that 
    $I\cap \bbF[X_{i_1},\dots,X_{i_d}]=(0)$. This is equivalent to say that 
    $[X_{i_1}],\dots, [X_{i_d}]\in A[V]$ are algebraicaly independent over $\bbF$. Therefore
    $d\leq \partial_{\bbF}(A[V])$. Now, suppose that $\varphi_i=[f_i]\in A[V],i=1,\dots m$ are 
    algebraicaly independent over $\bbF$ and put $N=\max{\deg(f_i)}$, then the map 
    \begin{equation}
        \bbF[Y_1,\dots, Y_m]_{\leq s}\mapsto \frac{\bbF[X_1,\dots,X_n]_{\leq Ns}}{I_{\leq Ns}}:g\mapsto g(\varphi_1,\dots,\varphi_m)
    \end{equation}
    is injective. Then, by taking dimensions we get
    \begin{equation}
     \binom{m+s}{s}  \leq \upleft{a}\text{HF}_{I}(s),\;\forall s\geq0. 
    \end{equation}
    We know that both expressions are polynomials in $s$, so by taking degrees we conclude that 
    $m\leq \dim(V) = d$ and $\partial_{\bbF}(A[V])\leq d$.\qed

    
    With small changes in the above proof, we can verify that in the case of $A[V]$ being an integral domain, we have 
    also $\partial_{\bbF}(K) = \dim(V),\;K=\text{Frac}(A[V])$.
\begin{example}
    Let $C=V(X_1-X_3^2,X_1-X_3^3)$, then $I(C)=(X_1-X_3^2,X_2-X_3^3)$ and 
    \begin{equation}
        \frac{\bbF[X_1,X_2,X_3]}{I(C)}\cong \bbF[X_3]
    \end{equation}
    so $\partial_{\bbF}(A[C])=1$.
\end{example}

\begin{example}
    The purpose of this example is to show that if $V\subset \bbF^n,|\bbF|=\infty$ is a linear space, then 
    $\dim(V)=\dim_{\bbF}(V)$.
    \begin{itemize}
        \item Let $\text{Ann}(V)=\text{Span}(f_1,\dots, f_m)$. Since $V$ has finite dimension, we have
        $v\in V\iff f(v)=0\;\forall f\in\text{Ann}(V)$ and $V=V(\{f_1,\dots,f_m\})$.
        \item Let $f_i=\sum a_{ij}X_j$, where we identify $X_j$ with the $j$th coordinate functions in $\bbF^n$. 
        Now, if $A=(a_{ij})$ and apply reduction to put $A$ in the form:
        \begin{equation}
            \begin{pmatrix}
                a_{1j_1}&*&0&*&\cdots&0&\cdots\\
                &&a_{2j_2}&*&\cdots&0&\cdots\\
                &&&&\vdots&\cdots&\vdots&\\
                &&&0&\cdots&a_{mj_m}&\cdots
            \end{pmatrix}
        \end{equation}
        with $a_{ij_i}=1,\;j_1<j_2<\cdots$. Now put LEX with: 
        $X_{j_1}>\cdots X_{j_m}>X_{j}\;j\notin\{j_1,\dots,j_m\}$ and write 
        $f_i = X_{j_i}+h_i$ with $h_i\in\bbF[X_j:j\notin\{j_1,\dots, j_m\}]$. Then we see that 
        $A[V]\cong\bbF[X_j:j\notin\{j_1,\dots, j_m\}]$ and $(f_1,\dots,f_m)$ is prime. Therefore 
        $I(V)=(f_1,\dots, f_m)$. Finally 
        \begin{equation}
            \dim(V)=\partial_{\bbF}(A[V]) = \dim_{\bbF}\bbF[X_j:j\notin\{j_1,\dots, j_m\}] = \dim_{\bbF}V.
        \end{equation}
    \end{itemize}
\end{example}
\end{itemize}


\section{Tangent Spaces}
\begin{itemize}
    \item(Regular maps).
    Let $V\subset\bbF^m, W\subset\bbF^n$ be algebraic sets. 
    Then a map $f:V\to W$ is said to be regular if there exist polynomials such that 
    $f=(f_1,\dots,f_m)$. Such a map is continuous in the Zariski topology. 
    \begin{remark}
        We can identify $A[V]$ with the set of regular functions $V\to\bbF$. In this 
        context, we have a contravariant functor:
        \begin{equation}
            \varphi: V\to W \rightsquigarrow \varphi^*:A[W]\to A[V]: [g]\mapsto [g]\circ\varphi
        \end{equation}
    \end{remark}
  \begin{example}
    Let $\varphi:\bbF\to C:x\mapsto (x^3,x^2,x)$ and write $C=(X_1-X_3^2,X_2-X_3^3)$. Remember the isomorphism
    $A[C]\to \bbF[Y]: [X_1]\mapsto Y^2,[X_2]\mapsto Y^3,[X_3]\mapsto Y$. So we have 
    \begin{equation}
        \varphi^*:A[C]\to A[\bbF]:[X_3]\mapsto [X_3]\circ\varphi. 
    \end{equation}
    The first map is $[X_3]:C\to \bbF:(x^3,x^2,x)\mapsto x$ and the last map is $\bbF\to C:x\mapsto x$. 
    At the algebraic level, the map $\varphi^*$ is just the the isomorphism $\bbF[Y]\cong\bbF[X]$.
  \end{example}
  \item(Tangent spaces). Let $V$ be an algebraic set and 
  $V_p = \bbF$ be the $A[V]-$module with action $[f]\alpha\equiv f(p)\alpha$. Let 
  $I(V) = (f_1,\dots, f_m)$. We define 
  \begin{equation}
    \begin{split}
        T_p(V) &= \{v\in\bbF^n: d(f)_p(v)=0\;\forall f\in I(V)\}\\
           &=\{v\in\bbF^n: d(f_i)_p(v)=0\;\forall i\}\\
           &=V(\{d(f_i)_p:i=1,\dots,m\})
    \end{split}
  \end{equation}

  \begin{remark}
    Caution ! The polynomials of $f_i$ here are not the generators of $V$, but of $I(V)$. For example, if 
    $V=V(f)$, we can conlude that $T_p(V) = V(df_p)$ only if $f$ is irreducible, i.e., if $I(V)=(f)$.
  \end{remark}
  Here $df_p\equiv \sum_{i=1}^nD_i(f)(p)X_i\in\bbF[X_1,\dots,X_n]$. We have an isomorphism
  \begin{equation}
    T_p(V)\rightarrow \text{Der}_{\bbF}(A[V],V_p):v\mapsto D_v;\;D_v([f])=\sum_{i=1}^n[D_i(f)]\cdot v_i\in\bbF.
  \end{equation}
  with inverse $D\mapsto (D([X_i]))_{i=1}^n$. Of course $D_v([f])=df_p(v)$.
  \begin{center}
    \textcolor{red}{Maybe it will be better to change notation of $V_p$ to $\upleft{p}\bbF$ or something.}
  \end{center}
  \begin{example}
    Let $C=(f_1,f_2),\;f_1=X_1-X_3^2, f_2=X_2-X_3^3$ and $p=(x^3,x^2,x)\in C$. We have 
    \begin{equation}
        d(f_1)_p = 1 - 2x X_3,\;d(f_2) = 1 - 3x^2 X_3 
    \end{equation}
    and $T_p(C)=\bbF\langle(3x^2,2x,1)\rangle$.
  \end{example} 

  \begin{example}
    The determinant is an irreducible polynomial. Indeed, let $\det\in\bbF[X_{ij}:1\leq i,j\leq n]$. Then, if we apply Laplace theorem
    on the first row, we get 
    \begin{equation}
        \det = A_{11} X_{11} - A_{12}X_{12} +\cdots+(-1)^nA_{1n}X_n
    \end{equation}
    where $A_{ij}$ denotes the determinant of the matrix obtained by excluding the $i$th row and $j$th collumn. By induction on 
    the number of variables, we assume that $A_{ij}$ are irreducible polynomials. Now, suppose that 
    $\det = fg$, where $\deg(f),\deg(g)>0$. Choose LEX with $X_{ij}>X_{k\ell}$ if $i<k$ or if $i=k$ and $j<\ell$. Then 
    \begin{equation}
     \LT(A_{11})X_{11} =  \LT(\det) = \LT(f)\LT(g).
    \end{equation}
    Assume that $X_{11}|\LT(f)$ (so $\LT(g)<X_{11}$) and write $f=AX_{11}+B,\;\LT(B)<X_{11}$. 
    Then we must have $AgX_{11}=A_{11}X_{11}$ and $A_{11}=Ag$. Since $\deg(g)>1$ and $A_{11}$ is irreducible, we must have 
    $g=A_{11}$ and $A=1$. Now we get 
    \begin{equation}
        - A_{12}X_{12} +\cdots+(-1)^nA_{1n}X_n = BA_{11}
    \end{equation}
    and $-\LT(A_{12})X_{12} = \LT(B)\LT(A_{11})$. Since $X_{22}|\LT(A_{11})$, we are led to $X_{22}|\LT(A_{12})$ which is 
    a contradiction because $X_{22}$ doenst appear in $A_{12}$ by definition.\par 

    Now, since $\det$ is an irreducible homogeneous polynomial, the difference $\det-1$ must be irreducible. Therefore
    $I(SL_n\bbF) = (\det - 1)$ and we can compute the tangent space $T_I(SL_n)=V(d_I(\det))$. 
    We have 
    \begin{equation}
        \begin{split}
            d_I\det &= \sum_{k,\ell} D_{k\ell}(\det)(I)X_{k\ell}\\
                &= \sum_{k,\ell}\sum_{\sigma\in S_n}D_{k\ell}\prod_{i=1}^nX_{i\sigma(i)}(I) X_{k\ell}\\
                &= \sum_{k,\ell}\sum_{\sigma\in S_n}\delta_{\ell,\sigma(k)}
                \prod_{i\neq k}\delta_{i,\sigma(i)} X_{k\ell}\\
                &= \sum_{k,\ell}\delta_{k\ell} X_{k\ell}\quad
                (\delta_{\ell,\sigma(k)}
                \prod_{i\neq k}\delta_{i,\sigma(i)}\neq0\;\text{only if }\sigma=1)\\
                &=\sum_k X_{kk}.
        \end{split}
    \end{equation}
    Therefore $T_ISL_n(\bbF) = \fr{sl}_n\bbF$, the Lie algebra of traceless matrices.
  \end{example}
  \item(Gauss lemma).
  We state the Gauss lemma.
  \begin{lemma}
    Let $R$ be a UFD ring (think in $\bbF[X]$), then $R[X]$ is a UFD.
  \end{lemma} 
  Proof: Let $f$ and $K=\text{Frac}(R)$. Since $K[X]$ is a UFD, we have 
  $f = FG,\;F,G\in K[X]$ with $F$ irreducible and $\deg(G)<\deg(f)$. Since $R$ is a UFD, there exist 
  $r,s\in R$ such that $rF,sG\in R[X]$. Decompose $rs = t_1\cdots t_k$ into irreducible elements. 
  Then if $F^\prime=rF, G^\prime=sG$, we have 
  \begin{equation}
    [F^\prime][G^\prime] = [rs][f] = 0 \in R/(t_i)[X]
  \end{equation}
  because $f\in R[X]$. So, since $R/(t_i)[X]$ is a domain, we see that $t_i|F^\prime$ or $t_i|G^\prime$. Then we see that 
  there exist a partition $\{1,\dots, k\}=I\cup J$ such that $F^\prime/\prod_{i\in I}t_i, G^\prime/\prod_{j\in J}t_j\in R[X]$.
  It's clear that $F^\prime$ still irreducible, then an induction on degree yields the result.\qed
  \begin{corollary}\label{corollary: generalized gcd}
    Let $f,g\in R=\bbF[X_1,\dots, X_n]$ and assume $f$ irreducible with $f\nmid g$, then there exist $F,G\in R$ such that 
    \begin{equation}
        Gf + Fg = h\in \bbF[X_1,\dots,X_{n-1}]\setminus0.
    \end{equation}
  \end{corollary}
  Proof: Let $K=\text{Frac}(\bbF[X_1,\dots,X_{n-1}])$, then $f,g\in K[X_n]$ and there exist 
  $A,B\in K[X_n]$ such that $Af + Bg = 1$ (by Gauss lemma $f$ still irreducible as an element of 
  $K[X_n]$). Then there exist $h\in \bbF[X_1,\dots,X_{n-1}]$ such that 
  $hA, hB\in \bbF[X_1,\dots,X_{n-1}]$. So by writing $G=hA, F =hB$ we are done.\qed
  \begin{notation}
    Let $A=\bbF[X_1,\dots,X_n]/I$, where $M$ is a $A-$module. We write 
    \begin{equation}
        \mathcal{T}_{A,M}=\{(v_1,\dots,v_n)\in M^n:\sum_i[D_i(f)]\cdot v_i,\;\forall f\in I\}.
    \end{equation}
    We have the isomorphism $\mathcal{T}_{A,M}\to \text{Der}_{\bbF}(A,M): v\mapsto D_v$ with inverse 
    \begin{equation}
        D\mapsto (D(\overline{X_i}))_{i=1}^n.
    \end{equation}
    Now consider the following context, let 
    $K=\text{Frac}(A)$ as $A$-module ($A$ acts by multiplication). Then we have an isomorphism
    \begin{equation}
        \text{Der}_{\bbF}(K,K)\to \text{Der}_{\bbF}(A,K):D\mapsto D|_A.
    \end{equation}
  \end{notation}
  In particular we also have the isomorphism $\mathcal{T}_{A,K} \cong  \text{Der}_{\bbF}(K,K)$. We observe that 
  this is also a $K-$vector space isomorphism, it's a bijection compatible with the $K-$action. From this isomorphism, 
  we see that $\mathcal{T}_{A,K}$ only depends on $K$, and the fact that $\text{Frac}(A)=K$.
  \begin{theorem}
    Let $A=\bbF[X_1,\dots,X_n]/I$ be a domain and $K=\text{Frac}(A)$. Then 
    \begin{equation}
        \partial_{\bbF}(A) = \dim_{K}\text{Der}_{\bbF}(K,K).
    \end{equation} 
  \end{theorem}
  Proof: Let $d=\partial_{\bbF}(A)=\partial_{\bbF}(K)$, then there exist 
  $z_1,\dots, z_{d+1}$ such that 
  \begin{itemize}
    \item $K=\bbF(z_1,\dots, z_{d+1})$
    \item $\{z_1,\dots,z_d\}$ is algebraicaly independent
    \item $z_d$ is separable over $\bbF(z_1,\dots, z_d)$
  \end{itemize}
  Let $R = \bbF[z_1,\dots, z_{d+1}]\subset K$ and $f\neq0$ irreducible such that 
  $f(z_1,\dots, z_{d+1})=0$. Then by corollary~(\ref{corollary: generalized gcd}), we see that 
  $\alpha:K[Y_1,\dots,Y_{d+1}]\to K:g\mapsto g(z_1,\dots, z_{d+1})$ has kernel $(f)$. 
  Therefore 
  \begin{equation}
    \begin{split}
        \mathcal{T}_{R,K} &= \{(v_1,\dots,v_{d+1})\in K^{d+1}:\sum_i[D_i(f)]\cdot v_i=0\}\\
        &=\{(v_1,\dots,v_{d+1})\in K^{d+1}: \sum_i D_i(f)(z_1,\dots, z_{d+1})\cdot v_i=0\}.
    \end{split}
  \end{equation}
  We see that $D_i(f)(z_1,\dots,z_{d+1})=0\Rightarrow D_i(f)\in (f)$ and $D_i(f)=0$. Since 
  $f\neq0$ and $|\bbF|=\infty$, we must have $D_i(f)\neq0$ for some $i$, so $D_i(f)(z_1,\dots, z_{d+1})\neq0$
  for some $i$. Then we conclude that 
  \begin{equation}
   \dim_K\text{Der}_{\bbF}(K,K)= \dim_{K} \mathcal{T}_{R,K} = d = \partial_{\bbF}(K) = \partial_{\bbF}(A).
  \end{equation}
  \qed
  \begin{corollary}
    If $A=\bbF[X_1,\dots,X_n]/(f_1,\dots,f_m)$, then  
    \begin{equation}
        \partial_{\bbF}(A) = \dim_K\mathcal{T}_{A,K}=n-\text{Rank}_K([D_i(f_j)]).
    \end{equation}
  \end{corollary} 

  \begin{example}
   Let $C=V(f_1,f_2), f_1 = X_1 - X_3^2, f_2 = X_2- X_3^3$. We have 
   \begin{equation}
    A[C] =\frac{\bbF[X_1,X_2,X_3]}{(f_1,f_2)}\to \bbF[X_3]:[g]\mapsto g(X_3^3,X_3^2,X_3).
   \end{equation}
   So $A[C]$ is a domain and we can define $K=\text{Frac}(A[C])$. We compute 
   \begin{equation}
    \text{Rank}_K([D_i(f_j)]) =
    \text{Rank}_K\begin{pmatrix}
       1&0& -2[X_3]\\ 0&1&-3[X_3]
    \end{pmatrix}=
    2.
   \end{equation}
   Then we conclude $\dim(V) = 3 - 2 =1$.
  \end{example}
  \item(differential of a regular map). Here we define the notion of differential of a regular map.
  \begin{definition}
    Let $\varphi:V\to W$ be a regular map between algebraic varieties. Also, let $p\in V, q = \varphi(p)\in W$. We define:
    \begin{equation}
        d\varphi_p: \text{Der}_{\bbF}(A[V],V_p)\to\text{Der}_{\bbF}(A[W],W_q):D_v\mapsto D_v\circ\varphi^*
    \end{equation}
  \end{definition}
  At the level of tangent spaces, we have:
  \begin{center}
    \begin{tikzcd}
        \text{Der}_{\bbF}(A[V],V_p)\arrow[rr,"d\varphi_p"]&&\text{Der}_{\bbF}(A[W],W_q)\arrow[dd]&\\
        &D_v\arrow[rr,mapsto,red]&&D_v\circ\varphi^*\arrow[dd, mapsto,red]\\
        T_p(V)\arrow[uu]\arrow[rr, dashed]&&T_q(W)&\\
        &v\arrow[uu,mapsto,red]\arrow[rr,mapsto,red]&&(D_v\circ\varphi^*([Y_i]))_{i=1}^m
    \end{tikzcd}
  \end{center}

  In more details, if we write $\varphi=(f_1,\dots, f_m)$, then $\varphi^*([Y_i])=[f_i]$ and we get 
  \begin{equation}
    D_v([f_i]) = \sum_{j=1}^m[D_j(f_i)]\cdot v_j = \sum_{j=1}^m D_j(f_i)(p)v_j.
  \end{equation}
  Therefore 
  \begin{equation}
    d\varphi_p(v) = (D_j(f_i)(p))_{i,j} \cdot v = J_p(\varphi)\cdot v.
  \end{equation}

  \begin{example}
    Let $G\subset \bbF^{n^2}$ be an algebraic group and $\mu:G\times G\to G$ be the product of matrices. 
    We have $G\times G\subset \bbF^{n^2}\times\bbF^{n^2}\cong \bbF^{2n^2}$. Now, identify the
    first $n^2$ coordinate functions with $X_{ij}$ and the second $n^2$ coordinate functions with $Y_{ij}$. So we have 
    \begin{equation}
        \mu_{ij} = \sum_\ell X_{i\ell}Y_{\ell j}.
    \end{equation} 
    Then we have 
    \begin{equation}
        \frac{\partial \mu_{ij}}{\partial X_{rs}}(I) = \frac{\partial \mu_{ij}}{\partial Y_{rs}}(I) =\delta_{ir}\delta_{sj}.
    \end{equation}
    Therefore 
    \begin{equation}
        d_{(I,I)}\mu(A,B) = (\sum_{r,s} \frac{\partial \mu_{ij}}{\partial X_{rs}}(I)A_{rs}+
          \frac{\partial \mu_{ij}}{\partial Y_{rs}}(I)B_{rs}  )_{ij} = (A_{ij} + B_{ij}) = A+B.
    \end{equation}
  \end{example}
  \begin{center}
    \TODO 
    \begin{itemize}
        \item[1.] Prove that $T_1G$ is a Lie algebra. Remember the star operator.
        \item[2.]  Compute the Lie algebra of $SL_n$ using the fact that $\det-1$ is an irreducible polynomial.
    \end{itemize}
  \end{center}
\end{itemize}
 

\section{Lecture 1 and 2}
First some definitions and lemmas worth to remember.
\begin{itemize}
    \item (Integral extension). Let $R$ be a ring and $S\subset R$ a subring. We say that 
    $R$ is integrally dependent over $S$ if every $r\in R$ is integral over $S$, that is, there 
    exist a monic polynomial $f\in S[X]$ such that $f(r)=0$.
    \begin{lemma}\label{lemma : characterization of integral extensions}
        Let $S\subset R$ be rings. 
        A finite subset of elements $b_1,\dots, b_m\in R$ are all integral over $S$ if and only if 
        $S[b_1,\dots, b_m]$ is finitely generated (as a $S-$module) over $S$.
    \end{lemma}
    Proof: $(\Rightarrow)$ Let $b\in R$ be integral over $S$, i.e., there exist 
    \begin{equation}
     f =    X^d + a_1X^{d-1} +\cdots + a_d \in S[X]
    \end{equation}
    such that $f(b)=0$. Then, using long division, we see that $S[b]=S +\cdots + S b^{d-1}$ which is finitely generated. 
    Now, if $b_1,\dots, b_m$ are all integral over $S$, we can argue by induction that $S[b_1,\dots, b_{m-1}]$ is finitely generated 
    over $S$. However, since $b_m$ is integral over $S[b_1,\dots, b_{m-1}],\;S[b_1,\dots,b_m]$ is finitely generated 
    over $[b_1,\dots, b_{m-1}]$ and it must be finitely generated over $S$ aswell.\par  
    $(\Leftarrow)$ Suppose that $S[b_1,\dots, b_m] = S\beta_1+\cdots+S\beta_n$ for some $\beta_i\in R$. Then for any 
    $b\in R$ which fixes $S[b_1,\dots,b_m]$, we have $b\beta_i= \sum_j a_{ij}\beta_j$ for some $(a_{ij})\in\text{Mat}_n(S)$. 
    Therefore 
    \begin{equation}
      0 =  (bI-(a_{ij}))^*(bI-(a_{ij}))v = \det(bI-(a_{ij}))v
    \end{equation}
    where $v=(\beta_1,\dots,\beta_n)^T$. Therefore $\det(bI-(a_{ij}))\beta_i=0,\forall i$ and since 
    $1\in S\beta_1+\cdots+S\beta_n$, we see that $\det(bI-(a_{ij}))=0$. Since $\det(X-(a_{ij}))\in S[X]$ is monic, we 
    are done.\qed 
    \begin{corollary}\label{corollary : transitive property integral dependence}
        (Transitive property of integral dependence). Given rings $R \supset S\supset T$, if $R$ is integral over $S$ and 
        $S$ is integral over $T$, then $R$ is integral over $T$.\par 
        \textcolor{red}{Question: Do we need Noetherian hypothesis using this statement?}
    \end{corollary}
    Proof: Let $r\in R$, then by hypothesis there exist $f=X^d+s_1X^{d-1}+\cdots\in S[X]$ such that $f(r)=0$.
    Define $A=T[s_1,\dots,s_d]$, then $A[r]$ is a finitely generated $A-$module because $S\subset A$ and lemma~(\ref{lemma : characterization of integral extensions}). 
    Now, since $A$ is a finitely generated $T-$module by lemma~(\ref{lemma : characterization of integral extensions}),
    we see that $A[r]$ is a finitely generated $T-$module. Now, \textcolor{red}{assuming $T$ is Noetherian} we see that 
    $T[r]$ is finitely generated as a $T-$module.\qed

    \begin{remark}
        We can certainly say (using the previous proof) that if $r\in R$, then $S[r]$ is finitely generated 
        as a $T-$module.
    \end{remark}

    \begin{lemma}\label{lemma : extension field transcendence degree}
        Let $K\supset k$ ($k$ a perfect field) be a finitely generated extension field with transcendence degree $\partial_k(K)=d$, then 
        there exist generators $z_1,\dots,z_d,z_{d+1}$ such that 
        \begin{itemize}
            \item[(1)] $K=k(z_1,\dots,z_{d+1})$
            \item[(2)] $\{z_1,\dots,z_d\}$ are algebraicaly independent
            \item[(3)] $z_{d+1}$ is separable over $k(z_1,\dots, z_d)$ 
        \end{itemize} 
    \end{lemma}  
\end{itemize}

\begin{dang}
\begin{lemma}
    (Noether's normalization lemma - Mumford's red book). Let $R$ be a finitely generated algebra over $k$. Assume also that $R$ is an integral domain.
    If $R$ has transcende degree $n$ over $k$ (which means that $\partial_k\text{Frac}(R)=n$), then 
    there exist $x_1,\dots, x_n\in R$ algebraicaly independent and such that $R$ is integral 
    over $k[x_1,\dots,x_n]$.
\end{lemma}
\end{dang}

\begin{dang}
    \begin{lemma}\label{lemma : Noether normalization}
        (Noether's normalization lemma - Statement made in class). Let $R$ be a finitely generated 
        $k-$algebra. Then there exist $x_1,\dots,x_n\in R$ algebraically independent over $k$ such that 
        $R$ is integral over $k[x_1,\dots,x_n]$.
    \end{lemma}
\end{dang}
Proof (Nagata): Since $R$ is finitely generated as a $k$-algebra, there exist generators $y_1,\dots,y_m$ such that 
$R= k[y_1,\dots, y_m]$. Now suppose that $y_1,\dots, y_m$ are algebraically independent, then we can take $x_i = y_i$. 
If not, then there exist $p\in k[Y_1,\dots, Y_m]\setminus0$ such that $p(y_1,\dots,y_m)=0$. Write 
\begin{equation}
    p(Y_1,\dots,Y_m)=\sum_{\alpha\in\bbZ^m_\geq0}a_\alpha Y_1^{\alpha_1}\cdots Y_m^{\alpha_m}
\end{equation}
and define a new polynomial 
\begin{equation}
    \begin{split}
        q(Y_1,\dots,Y_m) &= p(Y_1,Y_2+Y_1^{r_2},\dots, Y_m+Y_1^{r_m})\\
        &=\sum_\alpha a_\alpha(Y_1^{\alpha_1+r_2\alpha_2+\cdots+r_m\alpha_m}+\cdots) 
    \end{split}
\end{equation}
for some choice of positive integers $r_2,\dots,r_m$. Take a positive number $b>\alpha_i,\;\forall \alpha_i$ such that 
$a_\alpha\neq0$. If we take $r_i = b^i$, then $\alpha_1+\alpha_2b+\cdots +\alpha_mb^{m-1}$ are all different numbers in base $b$. 
So, from different powers of the form $\alpha_1+r_2\alpha_2+\cdots+r_m\alpha_m$ we have a maximal one 
$d=\max\{\alpha_1+\cdots r_m\alpha_m:a_\alpha\neq0\}$ and 
\begin{equation}
    q(Y_1,\dots,Y_m) = aY_1^d + \text{smaller degree terms} \in k[Y_2,\dots, Y_m][Y_1],\;a\neq0.
\end{equation}
Now, if we let $z_i=y_i - y_1^{r_i}\in R,i=2,\dots,m$, we see that $q(y_1,z_2,\dots,z_m)=p(y_1,\dots,y_m)=0$ and 
$y_1$ is integral over $k[z_2,\dots, z_m]$. By induction on number of generators, there exist $x_1,\dots, x_n$
algebraically independent such that $k[z_2,\dots, z_m]$ is integral over $k[x_1,\dots, x_n]$. So, by lemma~\ref{lemma : characterization of integral extensions}
and its corollary~(\ref{lemma : characterization of integral extensions}), we see that $k[y_1,z_2,\dots,z_m]$ is integral over $k[x_1,\dots,x_m]$. Now, by definition, it's clear 
that $k[y_1,z_2,\dots,z_m]$ is a finitely generated $R-$module. So by the corollary~(\ref{corollary : transitive property integral dependence}) again, we 
see that $R$ is integral over $k[x_1,\dots,x_m]$.\qed 

\begin{dang}
    \begin{lemma}\label{lemma : R is a field, S is a field}
        Let $S\subset R$ be a ring extension. If $R$ is a field, which integral over $S$, then $S$ is a field. 
    \end{lemma}
\end{dang}
Proof: Let $s\in S\setminus0$ and $r\in R$ be its inverse. Since $R$ is integral over $S$, there exist
$f=X^d+a_1X^{d-1}+\cdots\in S[X]$ be such that $f(r)=0$. Therefore 
\begin{equation}
   0= s^{d-1}f(r) = r + a_1 + sa_2+\cdots + s^{d-1}a_0\Rightarrow r\in S.
\end{equation} 
\qed

\begin{remark}
    Let $L|K$ be a field extension. Then
    \begin{center}
        \begin{tikzcd}
            L|K \text{ finite extension}\arrow[r, bend left= 30] &  L|K \text{ algebraic extension}
            \arrow[l, bend left= 30, "\text{+ finitely generated extension}"]
        \end{tikzcd}
    \end{center}
\end{remark}

\begin{dang}
    \begin{theorem}\label{theorem : Zariski 1}
        (Zariski). Let $k$ be a field and $\fr{m}\subset R=k[X_1,\dots, X_m]$ be a maximal ideal. Then $L=R/\fr{m}$ is 
        a finite extension field of $k$.
    \end{theorem}
\end{dang}
Proof: By Noether normalizaiton lemma~(\ref{lemma : Noether normalization}) there exist 
$a_1,\dots, a_n\in L$ algebraically independent such that $L$ is integral over $k[a_1,\dots,a_n]$. By lemma~(\ref{lemma : R is a field, S is a field}), 
we conclude that $k[a_1,\dots,a_n]$ is a field. Suppose that $n>0$, so $a_1$ has an inverse. This means that 
$a_1g(a_1,\dots, a_n) = 1$ for some $g\in k[X_1,\dots, X_n]$, which constradicts the fact that $a_1,\dots, a_n$ are 
alegrbaically independent. Then we conclude that $L$ is an integral (and then algebraic) over $k$. However, since 
$L$ is a finitely generated extension ($L=k(\overline{X_1},\dots,\overline{X_m})$) over $k$ we conclude that $L$ is in fact 
a finite field extension over $k$.\qed

\begin{dang}
    \begin{theorem}\label{theorem : Hilbert's zero theorem}
        Suppose that $k$ is an algebraically closed field.
        Let $I\subset R =k[X_1,\dots, X_m]$ be an ideal, then $V(I)=\varnothing\iff I = k[X_1,\dots, X_m]$.
    \end{theorem}
\end{dang}
Proof: If $I\neq R$, then there exist a maximal ideal $\fr{m}\supset I$. So $V(I)\supset V(\fr{m})$ and it's sufficient to 
prove that $V(\fr{m})\neq\varnothing$. By theorem~(\ref{theorem : Zariski 1}), we have that $L=R/\fr{m}$ is a finite 
$k-$extension, so we must have $L=k$ because $k$ is algebraically closed. So there exist $a_i\in k,i=1,\dots,m$ such that 
$X_i - a_i\in\fr{m}$ and $(X_1-a_1,\dots, X_m-a_m)\subset \fr{m}$. However $\varphi: R\to k:F\mapsto F(a_1,\dots,a_m)$ has kernel 
$\Ker(\varphi)=(X_1-a_1,\dots, X_m-a_m)$ and since $k$ is a field, $R/\Ker(\varphi)$ must be  a field and 
$\Ker(\varphi)$ is a maximal ideal. Therefore we conclude that $\fr{m}=(X_1-a_1,\dots, X_m-a_m)$. Finally
$V(X_1-a_1,\dots,X_m-a_m)=\{a_1,\dots,a_m\}\neq\varnothing$.\qed

\begin{remark}
    The hypothesis $k=\overline{k}$ in the previous theorem is crucial. For example if $k=\bbR$, we have 
    $V(x^2+1)=\varnothing$.
\end{remark}

\begin{remark}
    From the proof of the previous theorem, we see that there exist a bijection 
    \begin{equation}
        \begin{split}
            k^n &\longleftrightarrow \left\{\begin{matrix}\text{maximal ideals}\\\fr{m}\subset R\end{matrix}\right\}\\
           (a_1,\dots,a_n)&\longmapsto (X_1-a_1,\dots, X_n-a_n)
        \end{split}
    \end{equation}
\end{remark}